\section{Introduction}

Semantic entropy outperforms token entropy by 0.155 AUROC on TriviaQA --- then loses by 0.066 AUROC on TruthfulQA.
The same model (Llama-2-7B), the same uncertainty estimation method, two factual question-answering benchmarks, opposite orderings.
This is not noise: the reversal exposes a task-structure condition that determines when semantic-level uncertainty estimation helps and when it fails.

Knowing when a language model is likely to be wrong is a prerequisite for safe deployment.
Uncertainty estimation methods --- which score model outputs by predicted reliability --- are central to hallucination detection, selective abstention, and human-in-the-loop systems.
Yet practitioners selecting uncertainty methods face a fundamental question: which method should be used for which task?
Without benchmark-agnostic characterization, selection defaults to what performed best on the most popular evaluation dataset, with no guarantee of generalization.

The field has converged on four major uncertainty proxy types: token entropy (TE), which measures the spread of the model's output token distribution; semantic entropy (SE), which clusters semantically equivalent outputs before computing entropy; SelfCheckGPT consistency (SCG), which scores cross-sample agreement via BERTScore; and verbalized confidence (VC), where the model self-reports its certainty.
Prior work evaluates these methods in isolation or in pairs.
\citet{Kuhn2023} demonstrate SE outperforms TE at 65B scale on TriviaQA and NaturalQuestions --- but do not test at 7B scale, do not evaluate SCG or VC in the same study, do not measure the mechanism driving the gap, and do not test cross-benchmark generalization.
\citet{Xiong2023} systematically evaluate VC across scales, but without SE.
\citet{Huang2023} compare entropy-based and sampling-based methods without SE.
No prior study provides all four methods at 7B scale with mechanistic measurement and cross-task validation.

The gap matters precisely because 7B-scale models are the most widely deployed.
If the SE $>$ TE ordering established at 65B holds at 7B, practitioners have a principled basis for method selection.
If it does not --- or if the ordering is task-dependent --- benchmark-driven defaults may be systematically misleading.

Our key insight is that SE's advantage over TE is conditioned on output diversity structure: SE wins when incorrect model outputs are paraphrase-diverse (many different wrong answers that NLI clustering absorbs into distinct clusters), and loses when incorrect outputs are deterministically wrong (the same misconception repeated K times, all assigned to one cluster, yielding near-zero entropy despite consistent error).
Token entropy is indifferent to semantic equivalence --- it counts surface variation and therefore successfully flags deterministic-wrong outputs via low entropy.
On TriviaQA, incorrect answers are diverse; SE wins.
On TruthfulQA's adversarial misconceptions, incorrect answers are consistent; TE wins.

We measure this directly.
H-E1 establishes the SE $>$ TE gap (+0.155 AUROC, bootstrap CI on the gap excludes zero, N=98 TriviaQA).
H-M1 quantifies the mechanism: within NLI-equivalent paraphrase clusters, token entropy varies substantially (mean intra-cluster TE variance = 7.152 nats$^2$, 71$\times$ the 0.1 nats$^2$ threshold across 76/98 questions), confirming that TE aggregates surface-form noise that SE's clustering removes.
H-C1 tests generalization: on TruthfulQA, the ordering reverses (TE = 0.511 $>$ SE = 0.445), providing direct experimental evidence of the task-structure dependence.
H-M3 demonstrates that BERTScore-based SCG (AUROC = 0.378) is not equivalent to NLI-based SE (AUROC = 0.714) on short factual QA --- the two methods are mechanistically distinct.
H-M4 characterizes VC degeneracy at 7B scale (ECE = 0.430; 5 distinct confidence values; ${\sim}60\%$ of responses at 95\% confidence despite ${\sim}47\%$ empirical accuracy).

This paper makes four contributions:

\textbf{First,} we provide the first unified four-way comparison of SE, TE, SCG, and VC at 7B scale under identical experimental conditions on TriviaQA (N=98), confirming SE $>$ TE at sub-65B scale (gap = +0.155).

\textbf{Second,} we measure the paraphrase-noise mechanism: intra-cluster TE variance is 71$\times$ the practical significance threshold, directly demonstrating that TE's within-cluster surface variation is the driver of SE's discriminative advantage.

\textbf{Third,} we identify the task-structure scope condition: the SE $>$ TE advantage holds on open-domain factual recall (TriviaQA, where incorrect outputs are paraphrase-diverse) but reverses on adversarial misconception tasks (TruthfulQA, where incorrect outputs are deterministically wrong).
This scope condition was not identified in prior SE evaluations, which used only open-domain factual recall benchmarks.

\textbf{Fourth,} we characterize two failure modes of SE alternatives: BERTScore-based SCG fails on short QA answers (delta = 0.092 vs.\ 0.03 equivalence threshold) due to lexical-overlap limitations; VC produces degenerate near-constant confidence at 7B scale (ECE = 0.430).

Section~\ref{sec:related} reviews related work.
Section~\ref{sec:method} describes the methodology.
Section~\ref{sec:setup} presents the experimental setup.
Section~\ref{sec:results} reports results.
Section~\ref{sec:discussion} discusses implications and limitations.
Section~\ref{sec:conclusion} concludes.
